\documentclass[%singlesided,
               doublesided,
               paper=a4,
               fontsize=10pt
              ]{my-resume}
\NeedsTeXFormat{LaTeX2e}
\RequirePackage{xcolor}

%%%%%%%%%%%%%%%%%%%%%%%%%%%%%%%%%%%%%%%%%%%%%%%%%%%%%%%%%%%%%%%%%%%%%%%%%%%%%%%%
% set geometry
%%%%%%%%%%%%%%%%%%%%%%%%%%%%%%%%%%%%%%%%%%%%%%%%%%%%%%%%%%%%%%%%%%%%%%%%%%%%%%%%

\setlength\highlightwidth{8cm}
\setlength\headerheight{2cm}
\setlength\marginleft{1cm}
\setlength\marginright{\marginleft}     
\setlength\margintop{1cm}
\setlength\marginbottom{1cm}


%%%%%%%%%%%%%%%%%%%%%%%%%%%%%%%%%%%%%%%%%%%%%%%%%%%%%%%%%%%%%%%%%%%%%%%%%%%%%%%%
% FONTS
%%%%%%%%%%%%%%%%%%%%%%%%%%%%%%%%%%%%%%%%%%%%%%%%%%%%%%%%%%%%%%%%%%%%%%%%%%%%%%%%

\RequirePackage{fontspec}
\setmainfont{Roboto}

%%%%%%%%%%%%%%%%%%%%%%%%%%%%%%%%%%%%%%%%%%%%%%%%%%%%%%%%%%%%%%%%%%%%%%%%%%%%%%%%
% CONFIGURAÇÃO DOS LINKS (Remove as caixas coloridas)
%%%%%%%%%%%%%%%%%%%%%%%%%%%%%%%%%%%%%%%%%%%%%%%%%%%%%%%%%%%%%%%%%%%%%%%%%%%%%%%%
\usepackage{hyperref}
\hypersetup{
    colorlinks=true,      
    linkcolor=black,      
    filecolor=black,      
    urlcolor=black,       
    emailcolor=black
}

%%%%%%%%%%%%%%%%%%%%%%%%%%%%%%%%%%%%%%%%%%%%%%%%%%%%%%%%%%%%%%%%%%%%%%%%%%%%%%%%
% COLORS
%%%%%%%%%%%%%%%%%%%%%%%%%%%%%%%%%%%%%%%%%%%%%%%%%%%%%%%%%%%%%%%%%%%%%%%%%%%%%%%%

% Cores novas usando códigos Hexadecimais
\definecolor{deepblue}{HTML}{0B1D3A}     % Azul marinho bem escuro
\definecolor{lightgray}{HTML}{F0F4F8}    % Cinza azulado bem claro para o fundo
\definecolor{stargold}{HTML}{D4AF37}     % Dourado para os detalhes

% Associando as cores criadas aos elementos
\colorlet{highlightbarcolor}{lightgray}
\colorlet{headerbarcolor}{deepblue}

\colorlet{accent}{stargold}
\colorlet{heading}{deepblue}
\colorlet{emphasis}{deepblue}
\colorlet{body}{darkgray}

%%%%%%%%%%%%%%%%%%%%%%%%%%%%%%%%%%%%%%%%%%%%%%%%%%%%%%%%%%%%%%%%%%%%%%%%%%%%%%%%
% set document
%%%%%%%%%%%%%%%%%%%%%%%%%%%%%%%%%%%%%%%%%%%%%%%%%%%%%%%%%%%%%%%%%%%%%%%%%%%%%%%%

\begin{document}                

\begin{center}
    \color{white}
    {\Huge\bfseries Luísa Valias Ferreira} \\
    \vspace{0.15cm} % Pequeno espaço entre o nome e o subtítulo
    {\large Doutoranda em Astrofísica}
\end{center}
\vspace{1.5em}

\highlightbar{

    \section{Contato}
    
    \email{luisa24@ov.ufrj.br}
    \phone{+55 (41) 99779-4900}
    \linkedin{Luísa V. Ferreira}{https://www.linkedin.com/in/luísa-v-ferreira-120606}
    
    \section{Habilidades}
    
    \skillsection{Programação}
    \skill{Python}{5}
    \skill{LaTeX}{5}
    
    \vspace{0.5em}
    \skillsection{Sistemas Operacionais}
    \skill{Linux}{5}
    \skill{Windows}{4}
    
    \vspace{0.5em}
    \skillsection{Software \& Tools}
    \skill{Visualização}{5}
    (e.g. matplotlib, ...)\\
    \skill{Analise de dados}{5}
    (e.g. numpy, pandas, ...)\\
    \skill{Office}{4}
    
    \vspace{0.5em}
    \skillsection{Línguas}
    \skill{Inglês}{5}
    \skill{Espanhol}{5}
    \skill{Italiano}{3}
    \bigskip
    
    \section{Certificados}
    \simpleskill{Certificação de proficiência em inglês– 2024 Duolingo English Test - Nota: 135}
    \vspace{0.5em}
    \simpleskill{Curso de Astrofísica- 2022-2023- EdX Universidade Nacional da Austrália}

}
\mainbar{
    \section{Resumo}
    Mestre em Astronomia pela Universidade Federal do Paraná. Experiência no desenvolvimento de projetos de pesquisa internacionais, análise estatística avançada e modelagem computacional. Interessada nas áreas de Astrofísica, análise de dados e Astronomia Cultural.
    
    \section[\faGears]{EXPERIÊNCIAS PROFISSIONAIS}
    
    \textbf{2024 -- Presente} \\
    \textbf{Observatório do Valongo - UFRJ} \\
    \textit{Bolsista de Iniciação Científica}\\
    \vspace{-0.4cm}
    \begin{itemize}[leftmargin=0.4cm, topsep=0pt, itemsep=2pt]
        \item \small Desenvolver pesquisa focada em pequenos corpos do Sistema Solar.
        \item \small Analisar e organizar dados observacionais e imagens do asteroide troiano (1867) Deiphobus.
        \item \small Investigar características físicas do objeto por meio de dados de ocultação estelar para determinação de forma e tamanho.
    \end{itemize}
    
    \vspace{0.6em}

    \textbf{2024 -- 2025} \\
    \textbf{Observatório do Valongo - UFRJ} \\
    \textit{Extensionista no Valongo de Portas Abertas} \\
    \vspace{-0.4cm}
    \begin{itemize}[leftmargin=0.4cm, topsep=0pt, itemsep=2pt]
        \item \small Mediar e guiar visitações públicas e escolares às instalações do Observatório.
        \item \small Atuar na divulgação científica explicando conceitos de Astronomia, reconhecimento de constelações e história da ciência no Brasil.
        \item \small Operar e manusear telescópios durante sessões de observação noturna.
    \end{itemize}

    \vspace{0.6em}
    
    \textbf{2022 -- 2023} \\
    \textbf{Freelancer} \\
    \textit{Professora particular de inglês} \\
    \vspace{-0.4cm}
    \begin{itemize}[leftmargin=0.4cm, topsep=0pt, itemsep=2pt]
        \item \small Planejar e organizar cronogramas de aulas personalizadas e atividades didáticas focadas no desenvolvimento do estudante.
    \end{itemize}
    
    \vspace{0.6em}
    
    \textbf{2021 -- 2022} \\
    \textbf{Globalizando} \\
    \textit{Professora voluntária de inglês} \\
    \vspace{-0.4cm}
    \begin{itemize}[leftmargin=0.4cm, topsep=0pt, itemsep=2pt]
        \item \small Elaborar planos de aula e tarefas complementares de linguagem para projetos sociais de ensino.
    \end{itemize}

    \vspace{0.6em}
    
    \textbf{2021 -- 2022} \\
    \textbf{Universidade Federal do Paraná (UFPR)} \\
    \textit{Estagiária LAMIR-UFPR} \\
    \vspace{-0.4cm}
    \begin{itemize}[leftmargin=0.4cm, topsep=0pt, itemsep=2pt]
        \item \small Realizar a catalogação de dados e a organização de materiais e insumos de pesquisa no Laboratório de Análises de Minerais e Rochas.
        \item \small Criar conteúdos e posts para redes sociais com foco em divulgação científica, traduzindo rotinas laboratoriais e projetos de pesquisa para o público geral.
    \end{itemize}

}
\makebody
\clearpage


\pagestyle{highlightmain}
% The highlightbar needs to be filled to display mainbar contents correctly in singlesised mode
% For an empty highlightbar, fill with empty space
\highlightbar{\hfill}
\mainbar{

    
    \section[\faMortarBoard]{Educação}
    \textbf{2029 -- 2031}\\
    \textbf{Universidade Federal do Paraná (UFPR)\\
    \textit{Mestrado de Astronomia}

    \vspace{0.4em}
    
    \textbf{2024 -- 2029} \\
    \textbf{Universidade Federal do Rio de Janeiro (UFRJ)} \\
    \textit{Bacharelado em Astronomia}

    \section{Prêmios e Consquistas}
    \achievement{Menção honrosa na 14ª SIAC com o projeto "Estudo do troiano (1867) Deiphobus através de ocultações estelares" }

    \section{Soft Skills}
    \smallskip
    
    \tag{Análise de dados} \tag{Pesquisa Acadêmica} \\
    \medskip % Espaço vertical sutil entre as linhas de tags
    \tag{Divulgação Científica} \tag{Didática e Ensino} \\
    \medskip
    \tag{Organização de Tarefas} \tag{Proatividade}
    
}
\makebody

\end{document}
